\section*{Bab \ref{ch:g7:stats} --- Menata Data}

\begin{solution}{exo:g7:stats:1}
Membilang tiap nilainya:
\begin{center}
\renewcommand{\arraystretch}{1.2}
\begin{tabular}{l|ccccc|c}
nilai & $9$ & $11$ & $12$ & $14$ & $15$ & \omterm{def:g1:addition:def}{jumlah} \\
\hline
\omterm{def:g7:stats:frequency}{frekuensi} & $4$ & $3$ & $7$ & $3$ & $3$ & $20$ \\
\end{tabular}
\end{center}
\end{solution}

\begin{solution}{exo:g7:stats:2}
\omterm{def:g7:stats:frequency}{Frekuensi relatifnya}: $\frac{4}{20} = 20\,\%$, $\frac{3}{20} = 15\,\%$,
$\frac{7}{20} = 35\,\%$, $15\,\%$, $15\,\%$. \omterm{def:g1:addition:def}{Jumlahnya}:
$20 + 15 + 35 + 15 + 15 = 100\,\%$. \checkmark
\end{solution}

\begin{solution}{exo:g7:stats:3}
Survei pertama: $\frac{30}{50} = 0.6 = 60\,\%$. Survei kedua:
$\frac{45}{90} = 0.5 = 50\,\%$. Survei yang pertama menunjukkan
proporsi yang lebih tinggi, walaupun keluarga bermobil satunya lebih
sedikit dalam hitungan mutlak.
\end{solution}

\begin{solution}{exo:g7:stats:4}
$\intco{5}{15}$: $8, 12, 9, 14, 11$ --- \omterm{def:g7:stats:frequency}{frekuensinya} $5$.
$\intco{15}{25}$: $23, 22, 18, 16$ --- \omterm{def:g7:stats:frequency}{frekuensinya} $4$.
$\intco{25}{35}$: $31, 27, 29, 25, 33$ --- \omterm{def:g7:stats:frequency}{frekuensinya} $5$.
$\intco{35}{45}$: $35$ --- \omterm{def:g7:stats:frequency}{frekuensinya} $1$. \omterm{def:g1:addition:def}{Jumlahnya} $15$. \checkmark
\end{solution}

\begin{solution}{exo:g7:stats:5}
Empat batang setinggi $8$, $12$, $6$, $4$ di atas nilai $0$, $1$,
$2$, $3$ (skala setiap $2$ berjalan baik).
\end{solution}

\begin{solution}{exo:g7:stats:6}
\omterm{def:g1:addition:def}{Jumlahnya} $30$ keluarga. \omterm{def:g7:stats:frequency}{Frekuensi relatifnya}: $\frac{8}{30}$,
$\frac{12}{30}$, $\frac{6}{30}$, $\frac{4}{30}$. Sudutnya
($\times 360^\circ$): $96^\circ$, $144^\circ$, $72^\circ$, $48^\circ$;
\omterm{def:g1:addition:def}{jumlahnya} $96 + 144 + 72 + 48 = 360^\circ$. \checkmark
\end{solution}

\begin{solution}{exo:g7:stats:7}
Persennya: musim panas $\frac{162}{360} = 45\,\%$; musim semi
$\frac{90}{360} = 25\,\%$; musim gugur dan musim dingin
$\frac{54}{360} = 15\,\%$ masing-masing. Dari $60$ orang, musim panas
dipilih oleh $0.45 \times 60 = 27$ orang.
\end{solution}

\begin{solution}{exo:g7:stats:8}
Kelas A: $\frac{12}{30} = 40\,\%$. Kelas B: $\frac{14}{40} = 35\,\%$.
Kelas A punya proporsi pejalan kaki yang lebih tinggi, walaupun kelas
B punya pejalan kaki yang lebih banyak dalam hitungan.
\end{solution}

\begin{solution}{exo:g7:stats:9}
Dengan sumbu yang bermula dari $9\,700$, ketiga batangnya punya tinggi
yang terlihat $100$, $300$ dan $400$ satuan: yang terakhir tampak
empat kali yang pertama, seakan-akan penjualannya melipatempat. Dengan
sumbu dari $0$, batangnya hampir sama ($9\,800$ ke $10\,100$ adalah
kenaikan kira-kira $3\,\%$): gambar yang jujur menunjukkan penjualan
yang hampir datar.
\end{solution}

\begin{solution}{exo:g7:stats:10}
\omterm{def:g7:stats:frequency}{Frekuensi relatifnya} berjumlah $1$: $f = 1 - 0.35 - 0.4 = 0.25$.
\omterm{def:g7:stats:frequency}{Frekuensinya} (\omterm{def:g1:addition:def}{jumlahnya} $80$): $0.35 \times 80 = 28$;
$0.4 \times 80 = 32$; $0.25 \times 80 = 20$. Periksa: $28 + 32 + 20 = 80$.
\end{solution}

\begin{solution}{exo:g7:stats:11}
Murid perempuan: $55\,\%$ dari $400 = 220$; laki-laki: $180$. Murid
perempuan yang makan di kantin: $0.4 \times 220 = 88$; laki-laki:
$0.6 \times 180 = 108$. Seluruhnya: $88 + 108 = 196$ murid, yaitu
$\frac{196}{400} = 49\,\%$ dari sekolahnya --- di antara $40\,\%$ dan
$60\,\%$, lebih dekat ke angka murid perempuan karena mereka lebih banyak.
\end{solution}

\begin{solution}{pb:g7:stats:1}
\textbf{1.} \omterm{def:g7:stats:frequency}{Frekuensinya}: $120 > 90$, sekolah A menang. \omterm{def:g7:stats:frequency}{Frekuensi
relatifnya}: $\frac{120}{400} = 0.30 = 30\,\%$ melawan
$\frac{90}{250} = 0.36 = 36\,\%$: sekolah B menang. Hadiahnya
sebaiknya mengikuti \omterm{def:g7:stats:frequency}{frekuensi relatif} --- sekolah B mendaur ulang
bagian yang lebih besar dari yang dipakainya; sekolah A hanya minum
jus lebih banyak.

\textbf{2.} Guru: $\frac{18}{30} = 0.60 = 60\,\%$. Murid:
$\frac{210}{600} = 0.35 = 35\,\%$. Bersama-sama:
$\frac{18 + 210}{630} = \frac{228}{630} \approx 0.36 =
36\,\%$. Muridnya dua puluh kali lebih banyak, jadi \omterm{def:g7:stats:frequency}{frekuensi relatif}
gabungannya tertarik hampir sepenuhnya ke sisi mereka --- kelompok
yang besar membawa bobot yang besar.

\textbf{3.} \omterm{def:g7:stats:frequency}{Frekuensi relatif} harus berjumlah $100\,\%$
(\cref{def:g7:stats:frequency}), tetapi
$45 + 30 + 20 + 10 = 105$: sedikitnya satu juringnya salah.

\textbf{4.} \omterm{def:g7:stats:frequency}{Frekuensi} untuk \omterm{def:g7:stats:frequency}{frekuensi relatif} $0.15$:
$0.15 \times 40 = 6$. \omterm{def:g7:stats:frequency}{Frekuensi} sejauh ini: $14 + 10 + 6 = 30$, jadi
\omterm{def:g7:stats:frequency}{frekuensi} terakhirnya $40 - 30 = 10$, dengan \omterm{def:g7:stats:frequency}{frekuensi relatif}
$\frac{10}{40} = 0.25$. (Periksa: $0.35 + 0.25 + 0.15 + 0.25 =
1$.)

\textbf{5.} Sudutnya: $0.40 \times 360 = 144^\circ$,
$0.35 \times 360 = 126^\circ$, $0.25 \times 360 = 90^\circ$;
dan $144 + 126 + 90 = 360^\circ$: lingkarannya menutup.

\textbf{6.} $\frac{2.4}{10} = 0.24 = 24\,\%$ surat suaranya
kembali.

\textbf{7.} Buku telepon dan daftar mobil memuat sebagian besar rumah
tangga yang kaya, yang pilihannya berbeda dari pilihan negaranya;
$2.4$ juta jawaban itu adalah \omterm{def:g7:stats:frequency}{frekuensi} raksasa yang ditarik dari
penduduk yang salah. Seperti pada pertanyaan 1: yang penting bukan
berapa banyak yang kamu bilang, melainkan siapa --- \omterm{def:g7:stats:frequency}{frekuensi relatif}
yang dihitung pada sampel yang menyimpang menggambarkan sampelnya,
bukan negaranya.

\textbf{8.} Dukungan yang sebenarnya:
$0.7 \times 1\,000 + 0.3 \times 9\,000 = 700 + 2\,700 = 3\,400$
pendukung dari $10\,000$: yaitu $34\,\%$. Jajak pendapatnya:
$0.7 \times 500 + 0.3 \times 500 = 350 + 150 = 500$ dari
$1\,000$: yaitu $50\,\%$ --- enam belas angka terlalu tinggi, karena
yang kaya adalah setengah sampelnya tetapi sepersepuluh kotanya.

\textbf{9.} Jawabannya: $0.4 \times 200 = 80$ yang tidak puas dan
$0.1 \times 800 = 80$ yang puas, jadi $160$ jawaban yang $80$ di
antaranya tidak puas: ketidakpuasan yang terukur
$\frac{80}{160} = 50\,\%$. Ketidakpuasan yang sebenarnya:
$\frac{200}{1\,000} = 20\,\%$. Tidak ada yang berbohong --- yang
tidak puas memang lebih sering menjawab, dan surveinya mendengar
mereka lebih nyaring.

\textbf{10.} Proporsi yang sebenarnya: sepersepuluh kaya, jadi $100$
wawancara kaya dan $900$ wawancara sederhana:
$0.7 \times 100 + 0.3 \times 900 = 70 + 270 = 340$ dari
$1\,000$: jajak pendapat yang diperbaiki meramalkan $34\,\%$ --- yaitu
kebenaran pada pertanyaan 8. Lima puluh ribu wawancara yang dipilih
baik mengalahkan dua juta wawancara yang dipilih buruk.

\textbf{11.} Rumah sakit A: ringan $\frac{90}{100} = 90\,\%$; berat
$\frac{280}{400} = 70\,\%$; keseluruhan
$\frac{90 + 280}{500} = \frac{370}{500} = 74\,\%$.

\textbf{12.} Rumah sakit B: ringan $\frac{340}{400} = 85\,\%$;
berat $\frac{65}{100} = 65\,\%$; keseluruhan
$\frac{340 + 65}{500} = \frac{405}{500} = 81\,\%$. Rumah sakit A
menang pada kasus ringan ($90 > 85$) \emph{dan} pada kasus berat
($70 > 65$).

\textbf{13.} Secara keseluruhan, B menunjukkan $81\,\%$ melawan
$74\,\%$ milik A: rumah sakit yang kalah pada \emph{setiap} golongan
justru menang pada angka totalnya. (Pembalikan ini punya nama:
paradoks Simpson.)

\textbf{14.} Campurannya berlawanan: pasien A sebagian besar berat
($400$ dari $500$), pasien B sebagian besar ringan ($400$ dari $500$).
Kasus yang berat lebih jarang sembuh di mana pun ia dirawat, jadi
angka keseluruhan A terseret turun oleh kasus sulit yang diterimanya
--- angka keseluruhan adalah campuran berbobot, dan bobotnya berbeda
antara kedua rumah sakit itu. Pasien yang berat sebaiknya memilih
rumah sakit A: $70\,\%$ mengalahkan $65\,\%$ pada satu-satunya baris
yang menyangkut dirinya.

\textbf{15.} Misalnya: ``Perbandingan Anda mencampur pasien ringan
dan pasien berat, dan kedua rumah sakit itu merawat campuran yang
berlawanan --- kalau dipilah menurut keparahannya, rumah sakit A
menyembuhkan bagian yang lebih tinggi pada \emph{kedua} jenisnya.
Angka keseluruhan hanya boleh dibandingkan kalau kelompok di baliknya
serupa; kalau tidak, bobotnyalah yang menentukan pemenangnya, bukan
mutunya.''

\end{solution}
